\documentclass[titlepage]{article}

\usepackage[utf8]{inputenc}

\usepackage{graphicx}
\usepackage{caption}
\usepackage{subcaption}
\usepackage[normalem]{ulem}

\usepackage{pgfplots}
\pgfplotsset{compat=1.18} 

\title{
    \textbf{Traffic Simulator} \\ [0.4 em]
    \large User Manual
}
\begin{document}

\pagenumbering{gobble} 
\maketitle

\pagebreak

\tableofcontents

\pagebreak
\pagenumbering{arabic} 

\section{About}

This document describes how to use
Traffic Simulator %%
. The software was developed as part of the course
System- och programvaruutveckling (DT117G) at Örebro University.

\section{Simulation}

This section describes the everything related to the simulation mode, from
the inner workings of the simulation engine to its visual represation
and available user controls.

\subsection{Simulation engine}

Traffic Simulator %%
uses a simulation model based on the Nagel-Schreckenberger
traffic simulation model. The model uses discrete time, and hence
moves in textbf{steps} or textbf{ticks}.

\subsection{Simulation Visuals}

\begin{table}[h]
    \begin{center}
        \begin{tabular}{c p{0.8\textwidth}}
            \textit{car image} & A simulated car \\
            \textit{road image} & A piece of the road \\
            \textit{traffic light image} & Traffic light \\
        \end{tabular}
    \end{center}
\end{table}

\subsection{Maps} \label{maps}

The following maps are implemented*

\begin{itemize}
    \item \textbf{One way street}: A simple one way street. The first
        map made for this simulator.
    \item \textbf{T Intersection}: A one way street with a connecting
        one way street in the middle.
    \item \textbf{Highway}: Simple highway with two lanes in each
        direction. Overtakes are implemented, to allow faster driving
        vehicles to pass slower driving ones.
    \item \textbf{Autobahn}: A large highway with five lanes
        in each direction. Supports overtakes in all lanes, and
        outer lanes has an extra speed up, to ensure the overtaking
        car does not get stuck in the outer lane.
\end{itemize}

*more will be added in the future, as well as the possibility to create
custom maps. %%

\subsection{Live Statistics}

\begin{enumerate}
    \item \textbf{Simulated time} from when the simulation started.
    \item \textbf{Total cars spawned} during the simulation.
    \item The number of currently \textbf{stationary cars}.
    \item The current \textbf{number of cars} in the simulation.
    \item Current \textbf{average speed}.
    \item Current \textbf{inflow} and \textbf{outflow} of cars,
        respectively. Shown in vehicles per simulated minute.
\end{enumerate}

\subsection{Playback Controls}

\section{Settings Menu}

%% Insert image of the settings menu, with numbers for each setting

\begin{enumerate}
    \item \textbf{Map selector}
        Select the map, or scenario, to run a simulation on
        (see section \ref{maps} on page \pageref{maps}).
    \item \textbf{Road length}
        Set the number of cells in the main road.
        Only applicable to certain maps. Also note the scaling of vehicles
        may be very small for large values.
    \item \textbf{Offline simulation toggle}
        Allows for running simulations without visuals, and get the
        statistics summary at the end. Useful for large simulations since
        it greatly speeds up the simulation time.
    \item \textbf{Min speed}
        Set the minimum ideal speed on vehicles.
    \item \textbf{Max speed}
        Set the maximum ideal speed on vehicles.
    \item \textbf{Number of ticks}
        Only active when running offline simulations. This decides for how
        long the simulation will run.
    \item \textbf{Spawn rate}
        Set the chance of a vehicle spawning each tick.
    \item \textbf{Time scaling}
        Set the initial time scaling for the simulation. Not active for
        offline mode.
    \item \textbf{Start simulation}
        Start a new simulation with the current settings.
    \item \textbf{Exit program}
        Quit to desktop.
\end{enumerate}

\section{Statistics}

\end{document}
